\section{Arquitectura propuesta}
\label{sec:arquitectura}

\subsection{Vista de componentes}

PayCore es un monolito modular: un proceso, un despliegue y una base de datos.
La figura~\ref{fig:arquitectura} muestra la organización.

\begin{figure}[htbp]
\centering
\resizebox{0.95\textwidth}{!}{% Vista de componentes de PayCore.
\begin{tikzpicture}[
  font=\small,
  >={Stealth[length=2.2mm]},
  caja/.style={draw=grisborde, fill=white, rounded corners=3pt,
               minimum height=8mm, inner xsep=6pt, align=center},
  modulo/.style={caja, fill=acento!8, draw=acento!40, minimum width=21mm,
                 font=\scriptsize},
  flecha/.style={->, draw=tinta!65, thick},
]

% --- Fila de módulos de dominio (ancla de todo el bloque) ---
\node[modulo] (identity) at (0,0) {identity};
\node[modulo, right=2mm of identity] (accounts) {accounts};
\node[modulo, right=2mm of accounts] (transfers) {transfers};
\node[modulo, right=2mm of transfers] (ledger) {ledger};
\node[modulo, right=2mm of ledger] (idem) {idempotency};

\coordinate (centro) at ($(identity.west)!0.5!(idem.east)$);

% --- Capa HTTP, encima y centrada ---
\node[caja, minimum width=112mm, anchor=south]
  (routers) at ($(centro)+(0,10mm)$)
  {\codigo{api/}\ \ routers, esquemas Pydantic, autenticación\\[-3pt]
   {\scriptsize limitación de peticiones y manejo central de errores}};

% --- Capa de plataforma, debajo y centrada ---
\node[caja, minimum width=112mm, anchor=north] (platform) at ($(centro)-(0,10mm)$)
  {\codigo{platform/}\ \ configuración, unidad de trabajo, dinero\\[-3pt]
   {\scriptsize registro estructurado y seguridad}};

% --- Contenedor del backend ---
\begin{scope}[on background layer]
  \node[draw=grisborde, rounded corners=5pt, fill=grisclaro, inner sep=9pt,
        fit=(routers)(identity)(idem)(platform)] (api) {};
\end{scope}
% La etiqueta va a la izquierda del borde para no cruzarse con la flecha
% que entra por el centro.
\node[font=\scriptsize\bfseries, color=acento, anchor=south west]
  at ([shift={(2mm,1mm)}]api.north west) {Backend FastAPI};

% --- Cliente ---
\node[caja, minimum width=34mm, above=14mm of api.north]
  (ui) {React + Vite\\[-3pt]{\scriptsize TypeScript}};
\node[font=\scriptsize\bfseries, color=acento, above=1pt of ui] {Cliente};

% --- Base de datos ---
\node[caja, cylinder, shape border rotate=90, aspect=0.10,
      fill=acento!12, draw=acento!50, minimum width=62mm,
      below=14mm of api.south] (db)
  {PostgreSQL 16\\[-3pt]
   {\scriptsize users \textbullet\ accounts \textbullet\ transfers
    \textbullet\ ledger\_entries \textbullet\ idempotency\_keys}};

% --- Flechas principales ---
\draw[flecha] (ui.south) -- node[right, font=\scriptsize, xshift=1mm]
  {HTTPS, JSON, JWT} (ui.south|-api.north);
\draw[flecha] (api.south) -- node[right, font=\scriptsize, xshift=1mm]
  {una transacción por caso de uso} (db.north);

\draw[flecha] (routers.south) -- (transfers.north);


\end{tikzpicture}
}
\caption{Componentes de PayCore. El módulo \codigo{transfers} orquesta a
\codigo{accounts}, \codigo{ledger} e \codigo{idempotency} dentro de una única
transacción de base de datos gestionada por la unidad de trabajo de
\codigo{platform}.}
\label{fig:arquitectura}
\end{figure}

\subsection{Regla de frontera}

\textbf{Ningún módulo accede a las tablas de otro.} Las interacciones pasan
por los servicios de cada módulo. La regla es de disciplina, no está impuesta
por el compilador, y por eso queda documentada aquí y en el registro de
decisión correspondiente. \estado{implementado}

El núcleo de dominio no importa FastAPI ni conoce HTTP, lo que permite probar
la lógica de transferencia sin levantar un servidor. Se aplicó el estilo de
puertos y adaptadores solo donde se paga solo, sin convertirlo en una
ceremonia que multiplique las capas.

\subsection{Por qué no microservicios}

Repetido aquí porque es la pregunta más probable en una entrevista: la
transferencia debe ser atómica, y en un monolito con una base de datos la
atomicidad es un \codigo{COMMIT}. Separar cuentas y libro mayor en servicios
con bases de datos distintas obliga a implementar Saga o un commit en dos
fases, con compensaciones, estados intermedios observables y mucha más
superficie de error. Se estaría cambiando una garantía fuerte y gratuita por
una débil y cara, a cambio de una escalabilidad que este sistema no necesita.

Las fronteras entre módulos son, además, el mapa de una extracción futura si
alguna vez hiciera falta.

\subsection{Camino de evolución documentado}

La arquitectura orientada a eventos se descarta como núcleo pero se reconoce
como la vía natural para lo que ocurre \emph{después} de una transferencia:
notificaciones, puntuación de riesgo asíncrona y analítica. El libro mayor,
al ser de solo inserción, ya tiene la forma de un registro de eventos.
\estado{diseño}
